\cotefile{29}
\scanpage{77}
\mspage{\unsure{L|2|7}}

pour les modèles normaux (\unsure{résultant} de la partie facile de Lefschetz, en \unsure{vers.|version} pour les groupes fondamentaux rendus abéliens), on a fini. De façon plus précise il suffit, prenant un \inserted{modèle} $X$ \uline{projectif} sur $k$, (propre et normal), de prendre dans $X$ une ``courbe générique $C$'' et pour $\unsure{n|u}(K)$ le genre de la dite courbe générique.

% diagram in left margin, opposite the last two lines
\diagram{\begin{tikzcd}[ampersand replacement=\&,column sep=0.4em,row sep=1.1em]
 C' \arrow[d] \& {\subset} \& X' \arrow[d] \& \& \\
 C_{\unsure{k(u)|K(u)}} \& {\subset} \& X \& {\subset} \& P^{r}
\end{tikzcd}}

Y-a-t-il une démonstration du fait que cela donne une majoration, n'utilisant pas la théorie du groupe fondamental et le fait que $\underline{\mathrm{Pic}}^{\tau}_{X/k}$ est propre sur $k$ si $X$ est normale ?
% slanted
\margin{démonstration directe ???}
% archive sheet number: 76
